% Part: first-order-logic
% Chapter: sequent-calculus
% Section: rules-and-proofs

\documentclass[../../../include/open-logic-section]{subfiles}

\begin{document}

\olfileid{mvl}{seq}{rul}

\olsection{قاعدے تے \usetoken{P}{derivation}}

اگلی گل لئی منّ لو $\Gamma, \Delta, \Pi, \Lambda$،
!!{sentence}s دے محدود سلسلے نیں۔

\begin{defn}[سیکوئنٹ]
\emph{$n$ پاسیاں والا سیکوئنٹ} ایس روپ دی عبارت اے۔
\[
\Gamma_1 \nSequent \dots \nSequent \Gamma_n
\]
ایتھے ہر $\Gamma_1$ !!{sentence}s دا محدود سلسلہ اے، جیہڑا خالی وی ہو سکدا اے۔
ایہ جملے زبان~$\Lang L$ دے نیں۔
\end{defn}

\begin{defn}[ابتدائی سیکوئنٹ]
\emph{$n$ پاسیاں والا ابتدائی سیکوئنٹ} اوہ $n$ پاسیاں والا
سیکوئنٹ اے جیہدا روپ $!A \nSequent \dots \nSequent !A$ اے،
زبان دے کسے وی !!{sentence} $!A$ لئی۔

جے زبان وچ $0$-مقامی رابطہ~$\star$ اے، یعنی قضیاتی مستقل،
تاں اسیں سیکوئنٹ $\dots \nSequent
\star \nSequent \dots$ وی لیندے آں۔ ایس وچ $\star$ متعلقہ
صداقتی قدر~$\tf{\star} \in V$ والی تھاں اُتے آؤندا اے تے
باقی تھاواں خالیاں نیں۔
\end{defn}

$n$ قدراں والی منطق~$\Log{L}$ دے ہر رابطے لئی، ہر اوہ
صداقتی قدر دا اک منطقی قاعدہ اے جیہڑی اوہ رابطہ~$\Log{L}$
وچ لے سکدا اے۔ $n$ پاسیاں والے سیکوئنٹ حساب وچ~$\Log{L}$
دے !!^{derivation}s سیکوئنٹاں دے درخت نیں۔ سبھ توں اُتّے
ابتدائی سیکوئنٹ ہوندے نیں۔ جے کوئی سیکوئنٹ اک یا کئی ہوراں
دے ہِٹھ اے، تاں اوہ~$\Log{L}$ دے رابطیاں دے کسے استنباطی
قاعدے توں صحیح طرح نکلنا چاہیدا اے۔

\begin{defn}[قضیے]
!!^a{sentence}~$!A$، $n$ قدراں والی منطق~$\Log{L}$ دا
\emph{قضیہ} اے جے !!a{derivation} اوہ $n$ پاسیاں والے
سیکوئنٹ دا موجود ہووے جیہدے وچ $!A$ ہر اوہ تھاں اُتے آوے جیہڑی~$\Log{L}$ دی مقرر
صداقتی قدر نال جُڑی اے۔ اسیں $\Proves[\Log{L}]
!A$ لکھدے آں جے $!A$ قضیہ اے، تے جے ایہ قضیہ نہیں تاں $\Proves/[\Log{L}] !A$۔
\end{defn}

\begin{defn}[!!^{derivability}]
!!^a{sentence}~$!A$، !!{sentence}s دے سیٹ~$\Gamma$ توں
$n$ قدراں والی منطق~$\Log{L}$ وچ \emph{!!{derivable} اے}،
یعنی $\Gamma
\Proves[\Log{L}] !A$، اُدوں تے صرف اُدوں جدوں اک محدود
ذیلی سیٹ~$\Gamma_0 \subseteq
\Gamma$ تے اک سلسلہ $\Gamma_0'$ ہووے جیہڑا~$\Gamma_0$ دے
!!{sentence}s توں بنے، تے ہِٹھلے سیکوئنٹ دا
!!a{derivation} ہووے۔
\[ \Lambda_1 \nSequent \dots \nSequent \Lambda_n \]
ایتھے $\Lambda_i$، $!A$ اے جے تھاں $i$ مقرر صداقتی قدر نال
جُڑی ہووے، تے ہور صورت وچ $\Gamma_0'$ اے۔ جے $!A$،
$\Gamma$ توں !!{derivable} نہیں، تاں $\Gamma \Proves/ !A$
لکھدے آں۔
\end{defn}

مثلاً $3$ قدراں والی لوکاشیویچ دی منطق دا $3$ پاسیاں والا
سیکوئنٹ حساب اے۔ $3$ پاسیاں والے سیکوئنٹ $\Gamma \nSequent \Pi \nSequent
\Delta$ وچ $\Gamma$،~$\False$ نال، $\Delta$،~$\True$ نال،
تے $\Pi$،~$\Undef$ نال جُڑیا اے۔ مسلّمے $!A \nSequent !A \nSequent !A$
نیں۔ کیوں جے بس $\True$ مقرر اے، ایس لئی
$\Gamma \Proves[\LogLuk[3]] !A$ اُدوں تے صرف اُدوں اے جدوں
سیکوئنٹ $\Gamma \nSequent \Gamma \nSequent !A$ دا !!a{derivation}
ہووے۔ (جے $\Undef$ وی مقرر ہووے، تاں !!a{derivation} ہِٹھلے سیکوئنٹ دا چاہیدا ہووے گا۔
$\Gamma \nSequent !A \nSequent !A$۔)

\end{document}
